\section{Penanganan Error: Error Reporting dan Exception}

PHP memiliki sistem penanganan error yang fleksibel untuk mendeteksi, melaporkan, dan menangani masalah yang terjadi saat eksekusi. Tingkat error yang ditampilkan diatur melalui konfigurasi \code{error\_reporting} di \code{php.ini} atau secara runtime dengan fungsi \code{error\_reporting()}. Nilai umum: \code{E\_ALL} menampilkan semua error, \code{E\_ERROR} hanya error fatal, \code{E\_WARNING} untuk peringatan. Selama pengembangan (\textit{development}), aktifkan semua error: \code{error\_reporting(E\_ALL)} dan \code{ini\_set('display\_errors', 1)}. Di lingkungan produksi, matikan tampilan error dan arahkan ke file log untuk mencegah bocornya informasi sensitif \cite{phpManual}.

\begin{table}[ht]
\centering
\begin{tabular}{|P{3cm}|P{2cm}|P{7cm}|}
\hline
\textbf{Konstanta} & \textbf{Nilai} & \textbf{Deskripsi} \\
\hline
\code{E\_ERROR} & 1 & Error fatal; skrip berhenti \\
\hline
\code{E\_WARNING} & 2 & Peringatan; skrip tetap berjalan \\
\hline
\code{E\_NOTICE} & 8 & Pemberitahuan minor (variabel belum didefinisikan) \\
\hline
\code{E\_PARSE} & 4 & Error sintaks saat kompilasi \\
\hline
\code{E\_DEPRECATED} & 8192 & Fitur yang akan dihapus di versi mendatang \\
\hline
\code{E\_ALL} & --- & Semua error dan peringatan \\
\hline
\end{tabular}
\caption{Level error utama di PHP}
\end{table}

Sejak PHP 5, penanganan error modern menggunakan mekanisme \textbf{exception} dengan blok \code{try-catch-finally}. Exception adalah object yang merepresentasikan kesalahan yang terjadi saat runtime. Kode yang berpotensi gagal ditempatkan di blok \code{try}; jika terjadi exception, eksekusi melompat ke blok \code{catch} yang sesuai. Blok \code{finally} (opsional) selalu dijalankan, baik terjadi exception maupun tidak. PHP menyediakan beberapa class exception bawaan seperti \code{Exception}, \code{RuntimeException}, \code{InvalidArgumentException}, dan \code{TypeError} \cite{phpRightWay}.

\begin{webcode}[language=PHP, caption={Penanganan error dengan try-catch dan custom exception}]
<?php
// Custom exception
class ValidationException extends Exception {
  private array $errors;

  public function __construct(array $errors, int $code = 0) {
    $this->errors = $errors;
    parent::__construct("Validasi gagal", $code);
  }

  public function getErrors(): array {
    return $this->errors;
  }
}

function simpanData(string $nama, string $email): void {
  $errors = [];
  if (empty($nama))  $errors[] = "Nama wajib diisi.";
  if (!filter_var($email, FILTER_VALIDATE_EMAIL)) {
    $errors[] = "Format email tidak valid.";
  }
  if (!empty($errors)) {
    throw new ValidationException($errors);
  }
  echo "Data '{$nama}' berhasil disimpan.\n";
}

try {
  simpanData("", "bukan-email");
} catch (ValidationException $e) {
  echo $e->getMessage() . "\n";
  foreach ($e->getErrors() as $err) {
    echo "- {$err}\n";
  }
} catch (\Exception $e) {
  echo "Error umum: " . $e->getMessage();
} finally {
  echo "Proses selesai.\n";
}
?>
\end{webcode}

PHP juga menyediakan fungsi \code{set\_error\_handler()} untuk menangani error tradisional (E\_WARNING, E\_NOTICE) dengan cara khusus, dan \code{set\_exception\_handler()} untuk menangkap exception yang tidak di-catch. Fungsi \code{trigger\_error()} dapat digunakan untuk memicu error buatan. Pada PHP 8, banyak error yang sebelumnya hanya menghasilkan warning kini menghasilkan exception (misalnya, \code{TypeError} atau \code{ValueError}), membuat penanganan error menjadi lebih konsisten dan prediktabel. Praktik terbaik: selalu tangani error secara eksplisit, log error untuk analisis, dan berikan pesan error yang informatif namun tidak membocorkan detail implementasi \cite{phpManual}.

